%& -shell-escape
% !TeX root = thesis-text.tex
% !TeX encoding = UTF-8
% !TeX program = XeLaTeX

\documentclass{.local/lib/classes/ndvr-super-article-standalone}

\begin{document}


    \pagebreak

    \section{Подходы к определению понятия НДВ}


    Тема поиска нечётких дубликатов видео 
    относительно нова для~научного сообщества.
    И~потому существует целый ряд различных толкований и~определений.
    Наиболее употребляемые определения
    Ву \cite{Wu:2007a}, Шена \cite{Shen:2007},
    Башарата \cite{Basharat:2008}, и~Керубини \cite{Cherubini:2009},
    представленные далее, относятся к~одному и~тому~же.

    \subsection{Нечеткие дубликаты по~Ву}

    Нечеткие дубликаты видео~— идентичные или~приблизительно идентичные видео,
    почти точные копии друг друга,
    но~отличающийся
        форматами файлов,
        параметрами кодирования,
        светоизмерительными параметрами (цветность, яркость),
        применением монтажа (титры, логотип, рамка),
        длиной
        и~набором модификаций (вставка или~удаление кадров) \cite{Wu:2007a}.

    \begin{figure}[H]
        \includegraphics[width=\textwidth]%
            {.local/vec/out/pdf/ndv-missile-definition-1-wu}%
        \caption{
            Иллюстрация определения НДВ по~Ву \cite{Wu:2007a}.\\
            В~верхнем левом углу~— оригинал изображения.
        }
    \end{figure}


    \subsection{Нечеткие дубликаты по~Шену}

    Нечеткие дубликаты видео~— клипы,
    которые похожи или~почти дублируют друг от~друга, но~выглядят по-разному.
    Они~отличаются в~силу
        разных условий съемки
            (точка обзора и~настройки камеры, условия освещения,
            фон, передний план, и~т.~д.),
        преобразований
            (формат видео, частота кадров, разрешение, сдвиг, обрезка, гамма, контраст,
            яркость, насыщенность, размытие, выдержка, резкость и~т.~д.),
        операции монтажа
            (вставка, удаление, перемена мест кадров и~модификация содержимого)
        \cite{Shen:2007}.

    \begin{figure}[H]
        \includegraphics[width=\textwidth]%
            {.local/vec/out/pdf/ndv-missile-definition-2-shen}%
        \caption{
            Иллюстрация определения НДВ по~Шену \cite{Shen:2007}.\\
            В~верхнем левом углу~— оригинал изображения.
        }
    \end{figure}


    \subsection{Нечеткие дубликаты по~Башарату}

    Нечеткие дубликаты видео~— изображение того~же эпизода
    (например, «беспилотный летательный аппарат в~воздухе») при~различных точках обзора,
    размерах видео, ландшафте, погодных условиях, моделях БПЛА и~движениях камеры.
    Одна и~та~же смысловая концепция может
    быть выражена в~различной форме, при~разном освещении,
    и~параметрах сцены \cite{Basharat:2008}.


    \begin{figure}[H]
        \includegraphics[width=\textwidth]%
            {.local/vec/out/pdf/ndv-missile-definition-3-basharat}
        \caption{
            Иллюстрация определения НДВ по~Башарату \cite{Basharat:2008}.\\
            В~верхнем левом углу~— оригинал изображения.
        }
    \end{figure}

    \pagebreak

    \subsection{Нечеткие дубликаты по~Керубини}

    Нечеткие дубликаты видео~— приблизительно одинаковые видео,
    которые могут отличаться параметрами кодирования,
    светоизмерительными параметрами (цветность, яркость),
    применением монтажа (титры, логотип), или~наложением  аудио.
    Идентичные видео с~соответствующей дополнительной информацией,
    в~любом из~них (изменения длины клипа или~сцен) рассматриваются как~НДВ.
    Два различных видео с~разным представлением и~сценарии,
    считают НДВ, если они~используют одинаковую семантику,
    и~ни~одного из~них нет~дополнительной информации \cite{Cherubini:2009}.

    \begin{figure}[H]
        \includegraphics[width=\textwidth]%
            {.local/vec/out/pdf/ndv-missile-definition-4-cherubini}
        \caption{
            Иллюстрация определения НДВ по~Керубини \cite{Cherubini:2009}.\\
            В~верхнем левом углу~— оригинал изображения.
        }
    \end{figure}

    \pagebreak

    \subsection{Сравнение определений}

    На~данный момент до~сих пор не~существует единого
    определения нечетких дубликатов видео.
    Чаще всего в~работах используется определения Ву и~Шена.
    Однако определения сводятся, к~тому, что~НДВ должны нести очень
    близкую визуальную информацию, выражающих одинаковый смысл.

    Самое узкое определение дано в~работе Ву и~др.,
    Только идентичные и~приблизительно идентичные видео
    для~одного того~же сюжета с~незначительными фотографическими различиями
    и~правками расценены как~НДВ.
    Формально это определение совпадает
    с~понятием «копии видео» в~работе \cite{Law-To:2007}
    Видео-копией в~работе названы отрывки, идентичные по~смыслу,
    но~отличающиеся светоизмерительными или~геометрическими характеристиками.

    Аналогично проект TRECVID \cite{Over:2011} определяет копию, как
    «сегмент видео, полученный из~другого видео,
    как~правило, с~помощью различных преобразований,
    таких как~дополнение, удаление, модификация
    (ориентации, цвета, контраста, кодирование),
    перезапись и~т.~д.
    Для~задачи обнаружения копий на~основе 
    содержимого в TRECVID выделяют восемь
    трансформаций в~том числе:
        имитации аналоговой перезаписи,
        картинку в~картинке,
        вставку рисунка,
        повторное кодирование,
        изменение гаммы,
        снижения качества,
        компоновку видео
    и~комбинации из~трех случайно выбранных преобразований.

    В~работе Шена и~др. \cite{Shen:2007} определение НДВ распространяется
    на~фильмы с~одинаковой смысловой концепцией, но~снятые при~разных условиях.

    Башарат и~др. \cite{Basharat:2008}  
    предлагают еще~более общее определение.
    НДВ являются видео, изображающие одно и~то~же смысловое явление
    при~определенных обстоятельствах.


    В~целях вовлечения психофизического восприятия в~определение НДВ,
    Керубини и~соавторы \cite{Cherubini:2009} сделали масштабный интернет-опрос.
    Результаты этого опроса показывают, что~именно конечные пользователи
    считают нечеткими дубликатами.
    Человеческое восприятие НДВ отвечает
    многим свойствам формального определения
    относительно манипуляций с~визуальным содержанием
    \cite{Wu:2007a, Shen:2007}.

    Если в~сходные отрывки внесена дополнительная информация,
    то с точки зрения испытуемых, 
    они~уже не~воспринимаются как~нечеткие дубликаты.
    С~другой стороны, два разных видео с~различными героями
    при~различных сценариях оценены как~дубликаты, т.~к. они~имеют общий смысл,
    при~условии, что~ни~один из~роликов не~имеет дополнительной информации.

    Исследование подтверждает, что~пользователи считают 
    нечёткими дубликатами,
    видео, которые визуально подобны и~семантически идентичны.
    Это не~совпадает обычным поиском по~содержимому,
    т.~к. сходство контента является 
    единственным фактором, который нужно учитывать.
    А~в~случае с~пользовательским определением необходимо проверять
    и~семантику и~ее представление.

    Поиск НДВ может быть расценен как~мост
    между~традиционным поиском основанным на~подобии содержимого
    — у~видео должно быть подобное визуальное
    представление независимо от~семантики —
    и~смысловым видео-поиском
    — у~видео должна быть соответствующая
    семантика независимо от~представления.



    % Существует несколько различных понятий:
    % частичные нечеткие дубликаты, нечеткие дубликаты, копии, и~полные дубликаты.

    % Полные дубликаты семантически и~визуально идентичны с~исходным видео,
    % Они~имеют в~точности одинаковый сюжет, сценарий  и~т.~д.

    % Копии совместно используют ту~же семантику и~сцены что~и~оригинал,
    % но~отличаются визуальным представлением.
    % Копии могут быть получены в~результате 
    % разнообразных преобразований и~правок.
    % В~эту категорию попадают видео с~изменениями в~цвете,
    % яркостью, контрасте, частотой кадров, и~т.~д.,
    % или~с~добавлением баннеров, логотипов, и~т.~д.,
    % Таким образом, они~выражают в~точности одинаковые сюжет и~сцены,
    % но~с~визуальными различиями.

    % Нечеткие дубликаты используют одну и~ту~же семантику,
    % но~их внешний вид может отличаться.

    % Частичные нечеткие дубликаты, 
    % обычно имеют большое семантическое
    % и~визуальное несоответствие на~уровне целых видео.
    % Они~совместно используют похожие 
    % (нечетко дублирующихся) отрывки друг друга,
    % или~даже области отдельных изображений.
    % Например, два длинных видео с~совершенно разными сюжетами
    % но~со~сходными фрагментами или~областями включая
    % изображения в~одинаковых местах повествования.

\end{document}

